%!TEX root = main.tex
\section{Metodología}

\subsection{Descripción del Kernel de Suavizado}
El algoritmo implementado realiza un filtrado espacial de promedio (Media) sobre una ventana de $3 \times 3$ píxeles. El kernel recorre la imagen de entrada y calcula el valor de salida para cada píxel como la suma de sus 8 vecinos más él mismo, dividida por 9. Se implementó una lógica de "clamp" en los bordes para asegurar que los píxeles periféricos utilicen los valores de frontera, manteniendo la integridad de la resolución original.

\subsection{Plataforma Experimental}
Las pruebas se realizaron en un entorno Linux con acceso a herramientas de desarrollo de alto rendimiento:

\begin{itemize}
    \item \textbf{Procesador (CPU):} Arquitectura x86\_64 (24 hilos lógicos activos en las pruebas OpenMP).
    \item \textbf{Procesador (GPU):} NVIDIA GeForce RTX 3060 con 12 GB de VRAM.
    \item \textbf{Compiladores y Entornos:} 
        \begin{itemize}
            \item \texttt{g++ -O3} para las versiones Escalar y OpenMP (CPU).
            \item \texttt{nvcc} para la versión nativa de CUDA.
            \item \texttt{Python 3.10 + PyTorch 2.x} para la versión acelerada por librería.
        \end{itemize}
    \item \textbf{Imagen de Prueba (Heavy Stress):} Un conjunto de Julia generado sintéticamente con resolución 8K ($8192 \times 8192$ píxeles) en formato PGM P5 (Binario) para optimizar la carga de los 67 millones de puntos.
\end{itemize}

\subsection{Implementaciones Especiales}

\subsubsection{CUDA (Bajo Nivel)}
Se implementó un kernel nativo utilizando el modelo de memoria global. Se optimizaron las dimensiones de bloque a $16 \times 16$ hilos para maximizar la ocupación de los multiprocesadores de la GPU. El flujo sigue los "5 pasos mágicos": Malloc, Memcpy (H2D), Launch, Sync y Memcpy (D2H).

\subsubsection{PyTorch (Alto Nivel)}
Se desarrolló un script en Python que utiliza \texttt{torch.nn.functional.conv2d} con un kernel fijo de $1/9$. Esta versión permite evaluar el overhead de una librería generalista frente a una implementación específica.

\subsection{Procedimiento de Prueba}
Para cada versión, se realizaron 100 iteraciones acumulativas del suavizado (donde la salida de la iteración $n$ sirve como entrada para la $n+1$). Esta repetición es necesaria para amortizar los tiempos de inicialización y observar el comportamiento térmico y de ancho de banda del sistema bajo carga constante.

\FloatBarrier

\FloatBarrier